%==================================================================
% ABSTRAK (Bahasa Indonesia)
%
% Panduan singkat: satu paragraf, 150-250 kata, memuat berturut-turut
% (1) masalah, (2) tujuan, (3) metode, (4) hasil utama berupa angka,
% dan (5) simpulan. Abstrak TIDAK memuat sitasi, tabel, atau gambar.
% Pastikan setiap klaim di abstrak benar-benar ada buktinya di Bab 4.
%==================================================================
\phantomsection
\addcontentsline{toc}{chapter}{ABSTRAK}
\begin{center}
    \textbf{\judulid}\\[0.5cm]
    Oleh\\
    \penulis\\
    NIM: \nim\\[1.5em]
    \textbf{ABSTRAK}\\[0.5cm]
\end{center}

Sistem informasi akademik merupakan layanan yang digunakan seluruh mahasiswa,
namun keluhan mengenai kemudahan penggunaannya terus berulang setiap masa
pengisian rencana studi. Pengelola belum memiliki dasar yang terukur untuk
menentukan aspek mana yang perlu diperbaiki lebih dahulu. Penelitian ini
bertujuan mengukur kualitas layanan sistem informasi akademik, menganalisis
kesenjangan antara tingkat kepentingan dan tingkat kinerja, serta menentukan
prioritas perbaikannya. Pengukuran dilakukan menggunakan instrumen WebQual 4.0
yang terdiri atas 22 butir pada tiga dimensi, yaitu \textit{usability},
\textit{information quality}, dan \textit{service interaction quality}. Data
dikumpulkan melalui survei terhadap 96 mahasiswa aktif dan dianalisis secara
deskriptif, dilanjutkan dengan \textit{Importance-Performance Analysis} untuk
memetakan setiap butir ke dalam empat kuadran prioritas. Hasil penelitian
menunjukkan dimensi \textit{usability} memperoleh skor kinerja terendah,
sedangkan kesenjangan terbesar terjadi pada butir kemudahan navigasi. Pemetaan
kuadran menempatkan lima butir pada kuadran prioritas utama, seluruhnya
berkaitan dengan navigasi dan kejelasan alur pengisian. Penelitian ini
menyimpulkan bahwa perbaikan sebaiknya diarahkan pada arsitektur informasi
antarmuka, bukan pada penambahan fitur, dan menghasilkan daftar prioritas
perbaikan yang dapat langsung ditindaklanjuti pengelola.\\[0.6cm]

\noindent \textbf{Kata kunci:} \katakunci
\clearpage
